\documentclass[12pt, letterpaper]{article}
\usepackage[margin=1in]{geometry}
\usepackage{graphicx}
\usepackage{setspace}
\usepackage{amsmath}
\usepackage{amsfonts}
\usepackage{amssymb}
\usepackage{sectsty}
\usepackage{tikz}
\usepackage{xcolor}

\definecolor{antiquefuchsia}{rgb}{0.57, 0.36, 0.51}
\definecolor{carnationpink}{rgb}{1.0, 0.65, 0.79}
\usetikzlibrary{positioning}

\begin{document}

\begin{titlepage}
    \centering
    \begin{tikzpicture}[overlay, remember picture]
        % Background gradient
        \fill[black] (current page.south west) rectangle (current page.north east);
        \fill[left color=violet!80!black, right color=black!40!black] (current page.south west) rectangle (current page.north east);
        % Title
        \node[font=\sffamily\bfseries\Huge, text=white] (title) at (current page.center) {Calculus and Vectors};
        % Author
        \node[below=1cm of title, font=\sffamily\LARGE, text=white] {by Kensukeken};
        % Mathematical symbols and equations
        \node[below=2cm of title, font=\sffamily\Large, text=white] (equation) {$\vec{F} = m\vec{a}$};
        \node[below=1cm of equation, font=\sffamily\large, text=white] {Newton's Second Law};
        % ISBN or other info
        \node[below=1cm of current page.south, font=\sffamily\large, text=white] {ISBN: 123-456-789};
        \fill[antiquefuchsia] (title.east) ++(2cm,0) rectangle ++(5cm,-7cm);

    \end{tikzpicture}
\end{titlepage}

\clearpage % New page

\thispagestyle{empty} % No page number

\begin{center}
    \includegraphics[width=0.5\textwidth]{imgs/OIG4.png} % Adjust the width as needed
\end{center}

\vspace{2cm}

\begin{center}
    \Huge\textbf{Calculus and Vectors}
\end{center}

\vspace{2cm}

\begin{center}
    \Large{Kensukeken}
\end{center}

\end{document}
